\subsection{Hasil dan Perbandingan dengan \textit{Scikit-learn}}
\label{subsec:nb-hasil}

\todo{Bandingkan dengan \texttt{sklearn.naive\_bayes.GaussianNB} dan/atau \texttt{CategoricalNB} sesuai jenis fitur. Pakai \textit{preprocessing} dan \textit{fold} yang sama agar perbedaan hanya berasal dari implementasi \cite{pedregosa2011sklearn}.}

\begin{table}[H]
	\centering
	\caption{Perbandingan NB \textit{from scratch} dan \textit{scikit-learn}}
	\label{tab:nb-hasil}
	\begin{tabular}{l c c c c c}
		\toprule
		\textbf{Implementasi} & \textbf{Macro F1} & \textbf{Accuracy} & \textbf{Precision} & \textbf{Recall} & \textbf{Waktu (s)} \\
		\midrule
		\textit{From scratch} & & & & & \\
		\textit{Scikit-learn} & & & & & \\
		\bottomrule
	\end{tabular}
\end{table}

\begin{figure}[H]
	\centering
	\placeholdergambar{\textit{Confusion matrix} \textit{from scratch} vs \textit{scikit-learn}}
	\caption{\textit{Confusion matrix} NB}
	\label{fig:nb-cm}
\end{figure}


\subsubsection{Analisis}
\todo{Hasil seharusnya sangat dekat bila \textit{smoothing} sama. Jelaskan selisih akibat \texttt{var\_smoothing} atau penanganan kategori baru.}
